\documentclass[a4paper,10pt,twoside]{report}
\usepackage{geometry}\geometry{a4paper,top=3.5cm,bottom=3.5cm,%
left=2.5cm,right=2.5cm,heightrounded,bindingoffset=0mm}
\usepackage[T1]{fontenc}
\usepackage[utf8]{inputenc}
\usepackage[italian]{babel}
\usepackage{graphicx}
\usepackage[export]{adjustbox}%Per il Frame attrono le immagini e il valign
\usepackage{subfig}
\usepackage{amsmath,amsfonts,amssymb,braket,mathrsfs}
\usepackage{float}
\usepackage{tabularx,booktabs}
\usepackage{hyperref}
\usepackage{epsfig}
\usepackage{pdfpages} %Per gli allegati
%\usepackage{minipage}
\usepackage[output-decimal-marker={,}]{siunitx}
\DeclareSIUnit \days {gg}
\DeclareSIUnit\permille{\text{\textperthousand}} 
%ILe tre righe sotto ervono per mettere il grassetto dentro la tabelle siunitex usando \B, il rosso usando \RED, e il verde usando \GREEN
\sisetup{detect-weight,mode=text}
\usepackage{etoolbox}
\newrobustcmd\B{\DeclareFontSeriesDefault[rm]{bf}{b}\bfseries}
\newrobustcmd\RED{\DeclareFontSeriesDefault[rm]{bf}{b}\color{red}}
\newrobustcmd\GREEN{\DeclareFontSeriesDefault[rm]{bf}{b}\color{myGreen}}
%%
\usepackage{tikz}
\usepackage{pgfplots,pgfplotstable}
%\pgfplotsset{compat=1.15} %indica la versione da utilizzare per pgfplot
\usetikzlibrary{patterns} % per il tratteggio
\usepgfplotslibrary{groupplots}
\pgfplotsset{compat=newest}
%\usepackage{stanli}
\usepackage{xspace}% per lo spazio intelligente
\newcommand{\e}{\`E\xspace}  %E'
\usepackage{titlesec} % per formato custom dei titoli dei capitoli
%\usepackage{sideways}%%%
% redefinizione del formato del titolo del capitolo
      % da formato
      %   Capitolo X
      %   Titolo capitolo
      %   a formato
      %      Titolo capitolo 
	\titleformat{\chapter}
        {\normalfont\Huge\bfseries}{}{0em}{}
	\titlespacing*{\chapter}{0pt}{0in}{0.02in}
	\titlespacing*{\section}{0pt}{0.2in}{0.02in}
	\titlespacing*{\subsection}{0pt}{0.10in}{0.02in}
%serve per la didascalia di tabelle e figure:
\usepackage{caption}
\captionsetup{tableposition=top,figureposition=bottom,font=small}\captionsetup{format=hang,labelfont={bf,color=pantone186}} %didascalie a più righe allineate e il nome in grassetto
%non viene allineato a sinistra se la didascalia è corta una sola riga. PERCHé??
\usepackage{xcolor}
%serve per mettere il codice con lo sfondo grigio chiaro
\definecolor{pantone186}{RGB}{206, 17, 38} %il colore del logo UNITN
\definecolor{myGray}{gray}{0.5} %più basso più scuro è
\definecolor{myGreen}{rgb}{0.0, 0.5, 0.0}
\usepackage{listings} 
\lstset{basicstyle=\small\ttfamily,
backgroundcolor=\color{black!10},%
boxpos=c,%
stringstyle=\itshape,		
lineskip=3pt,%
numbers=left,
numberstyle=\footnotesize,}
\usepackage{lscape}
\usepackage{multirow}
\usepackage{import}

\usepackage{pifont}
\newcommand{\notchecked}{\text{\ding{55}}\xspace} 
\newcommand{\checked}{\text{\ding{51}}\xspace}


\begin{document}
\pagestyle{plain}
\thispagestyle{empty}
\begin{center}
  \begin{figure}[H]
    \centerline{\psfig{file=IMG/logo_unitn_black_centerNEW.eps,
                        width=0.8\textwidth,trim = 0 0.9cm 0 0.5cm}}
  \end{figure}
  \textcolor{pantone186}{\noindent\rule{\textwidth}{.5pt}}

  \Large\textsc{Dipartimento di Ingegneria Civile, Ambientale e Meccanica\\}
  \Large{Corso di Laurea in Ingegneria Civile
  }

  \vspace{3.2 cm} 
  %\Large\textsc{Elaborato finale\\} 
  %\vspace{1 cm} 
  \Huge\textsc{Relazione di Tecnica delle Costruzioni\\}
  
  \vspace{0.2 cm}
  \Large{\it{Progetto e verifica di alcuni elementi strutturali in C.A.\\
  di un edificio multipiano con struttura a telaio}}


  \vspace{4 cm} 
  \begin{tabular*}{\textwidth}{ l @{\extracolsep{\fill}} r }
  \Large\textsc{Docenti} & \Large\textsc{Studente}\\
  \Large{Nadia Baldassino}& \Large{Nicola Meoli 225077}\\
  \Large{Stefano Gasperetti}& \\

  	
  	
  \end{tabular*}

  \vspace{3.1cm} 
  \textcolor{pantone186}{\noindent\rule{\textwidth}{1pt}}
    
  \Large{Anno accademico 2021/22}
  
\end{center}

\tableofcontents
%\setcounter{page}{1}
%Tabelle e figure sulla stessa pagina:
%Le aggiunge all'indice. phantomsection serve per non far casini con hyperref
\clearpage
\begingroup
   %\let\cleardoublepage\relax  % book
    \let\clearpage\relax        % report
        \listoftables
        \phantomsection
        \addcontentsline{toc}{chapter}{Elenco delle tabelle}
\endgroup
        %
        \clearpage
\begingroup
        \let\clearpage\relax
        \listoffigures
        \phantomsection
        \addcontentsline{toc}{chapter}{Elenco delle figure}
\endgroup
%
%%%%%%%%%%
%Comandi aggiunti:
%%%%%%%%%%
\newcommand{\red}[1]{\textcolor{pantone186}{#1}}
\newcommand{\norma}[1]{\textsc{#1}}
\newcommand{\normaref}[1]{\texttt{#1}}
%

\chapter{Introduzione}
La presente relazione riguarda il progetto e la verifica strutturale di alcuni elementi componenti un edificio composto da 4 piani di cui 3 fuori terra sito in Provincia di Trento ad un'altitudine di \SI{300}{\meter} sopra il livello del mare.
Le dimensioni in pianta sono di circa $28 \times \SI{42}{\metre}$ per una altezza totale di \SI{12.45}{\meter} di cui \SI{9.70}{\meter} fuori terra. 
La cubatura totate è perciò di circa \SI{14650}{\metre\cubed}.

\paragraph{Composizione strutturale dell'edificio}
La struttura dell'edificio è a telaio in calcestruzzo armato, da progetto preliminare caratterizzata da pilastri quadrati di sezione \SI{30}{\centi\meter} e da travi sia a spessore che non, rispettivamente larghe \SI{60}{\centi\meter} ed alte come lo spessore del solaio, e $30 \times \SI{50}{\centi\metre}$.

I solai strutturali sono di tipo a lastre tralicciate Predalle tra il piano interrato e il piano terra e nel solaio di copertura. 
Il peso ultimato di tale solaio di \SI{3.6}{\kilo\newton\per\square\meter}.
I solai tra piano terra e piano primo e tra piano primo e piano secondo sono di tipo a travetti e latero-cemento.
Il peso ultimato di tale solaio di \SI{3.20}{\kilo\newton\per\square\meter}.

I solai dei piano intermedi sono costituiti da un pacchetto con calcestruzzo alleggerito, massetto di allettamento, intonaco e pavimento in ceramica.
Il solaio nella zona del terrazzo è costituito da isolante, impermealizzazione, massetto in calcestruzzo e pavimento.
Il solaio di copertura è costituito da isolante, massetto in calcestruzzo alleggerito, impermealizzazione e ghiaia.

Le pareti divisorie interne sono costituite da muratura in laterizio e intonaco in ambo i lati.
Le pareti perimetrali sono costituite da muratura in laterizio, cappotto esterno e intonaco.

Il piano terra è adibito a negozi, il piano primo ad uffici aperti al pubblico, il piano secondo a civile abitazione.
Il piano interrato è adibito a garage.

\paragraph{Oggetto della relazione}
Gli elementi che si andranno ad analizzare saranno una trave ed un solaio del primo piano e il pilastro P27. 
Tali elementi sono ben evidenziati nelle piante a seguire insieme alle sollecitazioni in essi agenti.

Subito dopo vengono elencati normative e materiali impiegati nell'analisi strutturale.

Seguono poi i capitoli di calcolo relativi agli elementi in esame.
La trave verrà analizzata per sollecitazioni agli SLU e SLE, solaio e pilastro per quelle agli SLU.

In allegato al presente documento sono riportati degli elaborati grafici eseguiti in software CAD.
\chapter{Piante e sollecitazioni}

\section{Azioni}
\clearpage	
\begin{landscape}
\subsection{SLU}
\begin{figure}[H]
\centering
\includegraphics[height=0.6\textwidth]{IMG/diagrammi_trave/m.pdf}
\caption{Diagrammi del momento generati con combinazione di carico SLU}
\label{fig:trave_ULS_momento}
\end{figure}
\begin{table}[H]
\footnotesize
\centering
\caption{Valori del momento con combinazione di carico SLU nei punti più significativi della struttura}
\label{tab:trave_ULS_momento}
	\begin{tabular}{lS[table-format=-3.2]S[table-format=-3.2]S[table-format=-3.2]S[table-format=-3.2]S[table-format=-3.2]S[table-format=-3.2]S[table-format=-3.2]S[table-format=-3.2]S[table-format=-3.2]S[table-format=-3.2]S[table-format=-3.2]S[table-format=-3.2]S[table-format=-3.2]}
		\toprule
		{} & {A1} & {C1} & {A2} & {C2} & {A3} & {C3} & {A4} & {C4} & {A5} & {C5} & {A6} & {C6} & {A7} \\
		\midrule
		$s\,\si{[\metre]}$ & 0.00 & 1.26 & 3.00 & 5.31 & 7.50 & 9.57 & 11.50 & 13.92 & 16.50 & 19.54 & 22.65 & 24.57 & 26.65 \\
        $M^{-}\,\si{[\kilo\newton\metre]}$ & 0.00 & -15.14 & -190.97 & 0.00 & -209.46 & -33.15 & -196.60 & -10.06 & -247.45 & 0.00 & -204.82 & -18.04 & -103.04 \\
        $M^{+}\,\si{[\kilo\newton\metre]}$ & 0.00 & 90.60 & 0.00 & 149.18 & 0.00 & 120.20 & 0.00 & 139.16 & 0.00 & 162.53 & 0.00 & 51.88 & 27.68 \\
		\bottomrule
	\end{tabular}
\end{table}
\end{landscape}

%{} & {A1} & {C1} & {A2} & {C2} & {A3} & {C3} & {A4} & {C4} & {A5} & {C5} & {A6} & {C6} & {A7} \\
%&\multicolumn{1}{c}{Nodo 1}&\multicolumn{1}{c}{Camp. 1}&\multicolumn{1}{c}{Nodo 2}&\multicolumn{1}{c}{Camp. 2}&\multicolumn{1}{c}{Nodo 3}&\multicolumn{1}{c}{Camp. 3}&\multicolumn{1}{c}{Nodo 4}&\multicolumn{1}{c}{Camp. 4}&\multicolumn{1}{c}{Nodo 5}&\multicolumn{1}{c}{Camp. 5}&\multicolumn{1}{c}{Nodo 6}&\multicolumn{1}{c}{Camp. 6}&\multicolumn{1}{c}{Nodo 7}\\
		

\chapter{Materiali}\label{chap:materiali}
\
\chapter{Trave: momento flettente SLU}
\section{Progetto della sezione maggiormente sollecitata}
In riferimento ai valori mostrati in precedenza nella tabella \ref{tab:trave_ULS_momento} a pagina \pageref{tab:trave_ULS_momento}, si progetta ora la sezione in calcestruzzo armato considerando quella maggiormente sollecitata.
Nel caso in esame è quella in appoggio 5 avente un momento pari a $\SI{-247.45}{\kilo\newton\metre}$.
Essendo soggetta a momento flettente negativo, sarà considerata invertita rispetto la dispozione reale e considerata quindi a momento flettente positivo, con le amrature invertite. 
A tale scopo verrà utilizzata la nomenclatura \emph{inv.} nelle tabelle a seguire per indicare tutte quelle sezioni che sono state inveritite.

\begin{figure}[ht]
  \centering
  \includegraphics[height=0.25\textheight]{IMG/IPE_slu_progetto.pdf}
  \caption{Convenzione e nomenclatura utilizzata per il progetto agli SLU}
  \label{fig:slu_progetto}
\end{figure}

Utilizzando la convenzione di segni riportata in figura \ref{fig:slu_progetto} e il polo $\mathbf{P}$ riportato, è possibile scrivere il sistema per determinare i valori dell'altezza utile $d$ e della quantita di armature inferiori $A_s$ che soddisfano il momento flettente applicato. 
\begin{equation}
  \begin{cases}
    C_c + C_s - T_s = 0 \\
    C_c \left(d - \lambda\,x\right) + C_s \left(d - d^{\prime\prime}\right) = M_{Ed}
  \end{cases}
\end{equation}
Sistema che è possibile risolvere esplicitando i valori delle risultanti e facendo delle ipotesi per ridurrre il numero di variabili. 
Si fissa il rapporto di armatura $\beta = 0.2$, ovvero pari al minimo possibile, in quanto il momente flettente a segno opposto nella sezione in esame è nullo.
Inoltre il progetto è vincolato, pertanto la base è imposta a \SI{300}{\milli\metre}.
Come copriferri $d^\prime,d^{\prime\prime}$ si scelgono entrambi uguali e pari a \SI{40}{\milli\metre}.
Infine si fissa il pieno utilizzo delle capacità meccaniche dei materiali: per cui l'acciaio superiore risulta snervato, e di conseguenza $\sigma_s^\prime = f_{yd}$, e il calcestruzzo a collasso, e di conseguenza $\sigma_c = f_{cd}$.
I coefficienti $\lambda$ e $\psi$ sono quelli del campo 3 e quindi pari a $\frac{17}{21}$ e $\frac{99}{238}$.

Il sistema espanso risulta:
\begin{equation}
  \begin{cases}
    b \, \psi \, \xi \, d \, f_{cd} + \beta \, A_s \, f_{yd} - A_s \, f_{yd} = 0 \\
    b \, \psi \, \xi \, d \, f_{cd} \left(d - \lambda\,x\right) + \beta \, A_s \, f_{yd} \left(d - d^{\prime\prime}\right) = M_{Ed}
  \end{cases}
\end{equation}
e risolvendolo si ottengono i seguenti valori:
\begin{equation}
  \hookrightarrow \quad
  \begin{cases}
    A_s &= \SI{1416.82}{\milli\metre\squared} \\
    A_s^\prime &= \SI{283.36}{\milli\metre\squared} \\
    d &= \SI{497.23}{\milli\metre}
  \end{cases}
\end{equation}

Si vuole cercare di avere un'altezza totale $h = \SI{500}{\milli\metre}$, per cui si sceglie un'altezza utile $d = \SI{460}{\milli\metre}$ inferiore a quella ottenuta dal sistema. 
I valori di armature $A_s$ sono soddisfatti con ad esempio 8Ø16, 6Ø18 o 5Ø20, aventi rispettivamente \SI{1608}, \SI{1527}, \SI{1570}{\milli\metre\squared} di area.
Avendo problemi riguardo la base minima su cui si dispongono le barre, il diametro da 16 non può essere utilizzato. 
Vengono perciò adottati i 6Ø18 che rispettano i requisiti, come di seguito dimostrato:
\begin{equation}  
  \begin{split}
    b_{min} &= 2 \cdot \text{copriferro} + 2 \cdot \varnothing_\textup{staffe} + n_\textup{barre} \cdot \varnothing_\textup{barre} + \left( n_\textup{barre} - 1 \right) \cdot \text{dist}_\textup{barre}  \\
    &= 2 \cdot 20 + 2 \cdot 10 + 6 \cdot 18 + \left( 6 - 1 \right) \cdot 25 \quad \si{\milli\metre}\\
    &= \SI{293}{\milli\metre} < b = \SI{300}{\milli\metre} \quad \checkmark ,
  \end{split}
\end{equation}
mentre per le armature superiori si scelgono 2Ø18 con unica funzione di reggistaffa.

\subsection{Verifica della sezione maggiormente sollecitata}
%--- LOG --- 
%Ipotesi di campo 3B, calcolo della retta limite 2-3
%xi_23 = 0.25926
%Esiste campo 2a, 2b, 3b perchè d2 = 40.00 < d2_xi23 = 55.77 mm
%xi_3 = 0.25166
%Ipotesi di campo 3 errata! xi_3 = 0.25166 < xi_23 = 0.25926 ma devo comunque verificare tramite es1 perché ho usato hp di es plasticizzato del 3B
%Essendo d2 = 40.00000 < d2_xi23 = 55.76720, non può esserci il campo 3A
%Calcolo xi con ipotesi di campo 2B: armature superiori snervate
%xi_2b = 0.25349
%!!!!!!!!!!!!!!verifica del ec da fare
%Ipotesi armature superiori snervate ok! es1 = 0.22308% > di ese1 = 0.18634%

%campo': '2B',
% 'ec': 0.0033956711163810735,
% 'es': 0.01,
% 'es1': 0.002230830149739241,
% 'xi': 0.2534901825283417,
% 'psi': 0.8036716031035525,
% 'lamb': 0.4138260346516767,
% 'Nrd': 0.0,
% 'Mrd': 247633547.16419485

  \begin{table}[htb]
    \centering
    \scriptsize
    \caption{ULS}
    \begin{tabular}{
        l
        c
        c
        S[table-format=3.3]
        l
        S[table-format=2.2]
        S[table-format=2.2]
        S[table-format=2.2]
        S[table-format=1.4]
        S[table-format=1.4]
        S[table-format=1.4]
        S[table-format=3.3]
        r
        S[table-format=2.1]}
    \toprule
    \multirow{2}{*}{Sez.}   & \multirow{2}{*}{$A_s$}    & \multirow{2}{*}{$A_s^\prime$} & {$M_{Ed}$}                    & \multirow{2}{*}{Campo}    & {$\varepsilon_c$}         & {$\varepsilon_s$}         & {$\varepsilon_s^\prime$}  & \multicolumn{1}{c}{\multirow{2}{*}{$\xi$}}    & \multicolumn{1}{c}{\multirow{2}{*}{$\psi$}}   & \multicolumn{1}{c}{\multirow{2}{*}{$\lambda$}}    & {$M_{Rd}$}                    & \multicolumn{2}{l}{$M_{Ed}<M_{Rd}$} \\
                            &                           &                               & {\si{[\kilo\newton\metre]}}   &                           &{\si{[\textperthousand]}}  &{\si{[\textperthousand]}}  &  {\si{[\textperthousand]}} &                        &                           &                               & {\si{[\kilo\newton\metre]}}   & & {\si{[\percent]}}\\
    \midrule
    A1 inv & 2Ø18 & 2Ø18 & 0.000   & 2A-1 & 1.51 & 10.00 & 0.51 & 0.1311 & 0.5648 & 0.3613 & 86.276  & \checkmark & 0.0 \\
    C1     & 3Ø18 & 2Ø18 & 90.595  & 2A-1 & 1.91 & 10.00 & 0.88 & 0.1607 & 0.6519 & 0.3724 & 128.015 & \checkmark & 70.8 \\
    A2 inv & 6Ø18 & 2Ø18 & 190.969 & 2B   & 3.40 & 10.00 & 2.23 & 0.2535 & 0.8037 & 0.4138 & 247.634 & \checkmark & 77.1 \\
    C2     & 4Ø18 & 2Ø18 & 149.180 & 2A-2 & 2.33 & 10.00 & 1.26 & 0.1890 & 0.7139 & 0.3856 & 168.988 & \checkmark & 88.3 \\
    A3 inv & 6Ø18 & 2Ø18 & 209.462 & 2B   & 3.40 & 10.00 & 2.23 & 0.2535 & 0.8037 & 0.4138 & 247.634 & \checkmark & 84.6 \\
    C3     & 3Ø18 & 2Ø18 & 120.204 & 2A-1 & 1.91 & 10.00 & 0.88 & 0.1607 & 0.6519 & 0.3724 & 128.015 & \checkmark & 93.9 \\
    A4 inv & 6Ø18 & 2Ø18 & 196.601 & 2B   & 3.40 & 10.00 & 2.23 & 0.2535 & 0.8037 & 0.4138 & 247.634 & \checkmark & 79.4 \\
    C4     & 4Ø18 & 2Ø18 & 139.165 & 2A-2 & 2.33 & 10.00 & 1.26 & 0.1890 & 0.7139 & 0.3856 & 168.988 & \checkmark & 82.4 \\
    A5 inv & 6Ø18 & 2Ø18 & 247.450 & 2B   & 3.40 & 10.00 & 2.23 & 0.2535 & 0.8037 & 0.4138 & 247.634 & \checkmark & 99.9 \\
    C5     & 4Ø18 & 2Ø18 & 162.530 & 2A-2 & 2.33 & 10.00 & 1.26 & 0.1890 & 0.7139 & 0.3856 & 168.988 & \checkmark & 96.2 \\
    A6 inv & 6Ø18 & 2Ø18 & 204.824 & 2B   & 3.40 & 10.00 & 2.23 & 0.2535 & 0.8037 & 0.4138 & 247.634 & \checkmark & 82.7 \\
    C6     & 2Ø18 & 2Ø18 & 51.884  & 2A-1 & 1.51 & 10.00 & 0.51 & 0.1311 & 0.5648 & 0.3613 & 86.276  & \checkmark & 60.1 \\
    A7 inv & 3Ø18 & 2Ø18 & 103.039 & 2A-1 & 1.91 & 10.00 & 0.88 & 0.1607 & 0.6519 & 0.3724 & 128.015 & \checkmark & 80.5 \\
    \midrule
    A1     & 2Ø18 & 2Ø18 & 0.000   & 2A-1 & 1.51 & 10.00 & 0.51 & 0.1311 & 0.5648 & 0.3613 & 86.276  & \checkmark & 0.0 \\
    C1 inv & 2Ø18 & 3Ø18 & 15.137  & 2A-1 & 1.42 & 10.00 & 0.42 & 0.1241 & 0.5410 & 0.3591 & 86.200  & \checkmark & 17.6 \\
    A2     & 2Ø18 & 6Ø18 & 0.000   & 2A-1 & 1.26 & 10.00 & 0.28 & 0.1119 & 0.4978 & 0.3555 & 86.004  & \checkmark & 0.0 \\
    C2 inv & 2Ø18 & 4Ø18 & 0.000   & 2A-1 & 1.35 & 10.00 & 0.36 & 0.1189 & 0.5230 & 0.3575 & 86.126  & \checkmark & 0.0 \\
    A3     & 2Ø18 & 6Ø18 & 0.000   & 2A-1 & 1.26 & 10.00 & 0.28 & 0.1119 & 0.4978 & 0.3555 & 86.004  & \checkmark & 0.0 \\
    C3 inv & 2Ø18 & 3Ø18 & 33.154  & 2A-1 & 1.42 & 10.00 & 0.42 & 0.1241 & 0.5410 & 0.3591 & 86.200  & \checkmark & 38.5 \\
    A4     & 2Ø18 & 6Ø18 & 0.000   & 2A-1 & 1.26 & 10.00 & 0.28 & 0.1119 & 0.4978 & 0.3555 & 86.004  & \checkmark & 0.0 \\
    C4 inv & 2Ø18 & 4Ø18 & 10.062  & 2A-1 & 1.35 & 10.00 & 0.36 & 0.1189 & 0.5230 & 0.3575 & 86.126  & \checkmark & 11.7 \\
    A5     & 2Ø18 & 6Ø18 & 0.000   & 2A-1 & 1.26 & 10.00 & 0.28 & 0.1119 & 0.4978 & 0.3555 & 86.004  & \checkmark & 0.0 \\
    C5 inv & 2Ø18 & 4Ø18 & 0.000   & 2A-1 & 1.35 & 10.00 & 0.36 & 0.1189 & 0.5230 & 0.3575 & 86.126  & \checkmark & 0.0 \\
    A6     & 2Ø18 & 6Ø18 & 0.000   & 2A-1 & 1.26 & 10.00 & 0.28 & 0.1119 & 0.4978 & 0.3555 & 86.004  & \checkmark & 0.0 \\
    C6 inv & 2Ø18 & 2Ø18 & 18.039  & 2A-1 & 1.51 & 10.00 & 0.51 & 0.1311 & 0.5648 & 0.3613 & 86.276  & \checkmark & 20.9 \\
    A7     & 2Ø18 & 3Ø18 & 27.682  & 2A-1 & 1.42 & 10.00 & 0.42 & 0.1241 & 0.5410 & 0.3591 & 86.200  & \checkmark & 32.1 \\
    \bottomrule
    \end{tabular}
    \end{table}
\chapter{Trave: Taglio SLU}
Si progettano ora le armature trasversali resistenti a taglio. 
Si sceglie di usare staffe ovunque e non ferri piegati, pertanto l'inclinazione di esse sarà sempre pari ad $\alpha = \SI{90}{\degree}$.

Agli appoggi, dove quindi il taglio è maggiore, si verificherà prima se è sufficiente la resistenza a taglio senza armatura specifica, in caso contrario si procederà al progetto della stessa.
In campata invece, dove il taglio è pressoché nullo, si utilizzerà sicuramente l'armatura trasversale minima.

\section{Armatura a taglio minima}
Le \norma{NTC18} fissano l'obbligo di disporre un quantitativo minimo di armatura trasversale indipendentemente dalla resistenza della trave. 
Tale quantitativo è calcolabile come:
\begin{equation}
    A_{st, min} = \num{1.5} \, b \si{\milli\metre\squared\per\metre} = \SI{450}{\milli\metre\squared\per\metre} \quad ,
\end{equation}
con un passo:
\begin{equation}
    s_{max} = \min \left\{ 0.8\,d ; \frac{\SI{1000}{\milli\metre}}{3} \; 15\, \varnothing_{long} \right\} = \SI{270}{\milli\metre}
\end{equation} 

\section{Progetto e verifica agli appoggi}
\paragraph{Verifica senza armatura specifica a taglio}
\begin{equation}
    V_{Rd} = \max \left\{ \left(\frac{0.18 \, k \, (100 \, \rho_l \, f_{ck})^{\tfrac{1}{3}}}{\gamma_c} + 0.15 \, \sigma _{cp}\right) b_w d \; ; \;(\nu_{min} + 0.15\sigma _{cp}) \, b_w \, d \right\}
\end{equation}
dove 
\[
    \begin{split}
        k &= \min \left\{ 1 + \left(\frac{200}{d}\right)^ {\tfrac{1}{2}} \; ; \; 2\right\} \\
        \nu_{min} &= 0.035 k ^ \frac{3}{2} f_{ck}^ \frac{1}{2} \\
        \rho_l &= \min \left\{ \frac{A_{sl}}{b_w d} \; ; \; 0.02\right\} \\
        \sigma_{cp} &= \min \left\{ \frac{N_{Ed}}{A_c} \; ; \; 0.2 f_{cd} \right\} \\
    \end{split}
\]
e in cui $A_{sl}$ corrisponde all'area tesa longitudinale $A_s$ e $b_w$ coincide con la base $B$ della sezione rettangolare.

Tale procedimento è stato applicato per tutti gli appoggi e come si evince dalla tabella \ref{tab:V_no_specific_armor} non è sufficiente a garantire la resistenza a taglio. 
\begin{table}[htb]
    \centering
    \scriptsize
    \caption{ULS taglio senza armatura}
    \label{tab:V_no_specific_armor}
    \begin{tabular}{
        l
        c
        S[table-format=3.3]
        S[table-format=1.2]
        S[table-format=1.2]
        S[table-format=1.2]
        S[table-format=1.2]
        S[table-format=3.3]
        S[table-format=3.3]
        S[table-format=3.3]
        c}
    \toprule
    \multirow{2}{*}{Sez.} & \multirow{2}{*}{$A_{sl}$} & {$V_{Ed}$} & {\multirow{2}{*}{$k$}} & {\multirow{2}{*}{$\nu_{min}$}} & {$\rho_l$} & {$\sigma_{cp}$} & {$V_{Rd,1}$} & {$V_{Rd,2}$} & {$V_{Rd}$} &  \multirow{2}{*}{$V_{Ed}<V_{Rd}$} \\
     & &{\si{[\kilo\newton]}} & & &{\si{[\percent]}}&{\si{[\mega\pascal]}}&{\si{[\kilo\newton]}}&{\si{[\kilo\newton]}}&{\si{[\kilo\newton]}}& \\
    \midrule
    A1 sx & 2Ø18 & 0.000   & 1.66 & 0.37 & 0.37 & 0.00 & 57.621 & 0.000 & 57.621 & \checked \\
    A1 dx & 2Ø18 & 144.220 & 1.66 & 0.37 & 0.37 & 0.00 & 57.621 & 0.000 & 57.621 & \notchecked \\
    A2 sx & 2Ø18 & 236.080 & 1.66 & 0.37 & 0.37 & 0.00 & 57.621 & 0.000 & 57.621 & \notchecked \\
    A2 dx & 6Ø18 & 275.030 & 1.66 & 0.37 & 1.11 & 0.00 & 83.103 & 0.000 & 83.103 & \notchecked \\
    A3 sx & 2Ø18 & 275.730 & 1.66 & 0.37 & 0.37 & 0.00 & 57.621 & 0.000 & 57.621 & \notchecked \\
    A3 dx & 6Ø18 & 263.720 & 1.66 & 0.37 & 1.11 & 0.00 & 83.103 & 0.000 & 83.103 & \notchecked \\
    A4 sx & 2Ø18 & 254.150 & 1.66 & 0.37 & 0.37 & 0.00 & 57.621 & 0.000 & 57.621 & \notchecked \\
    A4 dx & 6Ø18 & 236.150 & 1.66 & 0.37 & 1.11 & 0.00 & 83.103 & 0.000 & 83.103 & \notchecked \\
    A5 sx & 2Ø18 & 255.130 & 1.66 & 0.37 & 0.37 & 0.00 & 57.621 & 0.000 & 57.621 & \notchecked \\
    A5 dx & 6Ø18 & 235.450 & 1.66 & 0.37 & 1.11 & 0.00 & 83.103 & 0.000 & 83.103 & \notchecked \\
    A6 sx & 2Ø18 & 227.870 & 1.66 & 0.37 & 0.37 & 0.00 & 57.621 & 0.000 & 57.621 & \notchecked \\
    A6 dx & 6Ø18 & 184.800 & 1.66 & 0.37 & 1.11 & 0.00 & 83.103 & 0.000 & 83.103 & \notchecked \\
    A7 sx & 2Ø18 & 149.290 & 1.66 & 0.37 & 0.37 & 0.00 & 57.621 & 0.000 & 57.621 & \notchecked \\
    A7 dx & 3Ø18 & 0.000   & 1.66 & 0.37 & 0.55 & 0.00 & 65.959 & 0.000 & 65.959 & \checked \\
    \bottomrule
    \end{tabular}
    \end{table}
    



\paragraph{Progetto del passo staffe}
Occorre perciò progettare un'armatura specifica che resista alle sollecitazioni taglianti.
In particolare si fa riferimento al capitolo \normaref{4.1.2.3.5.2} delle \norma{NTC18}, il quale specifica che la resistenza di progetto a taglio è la minore tra la resistenza a taglio trazione $V_{Rsd}$ e la resistenza taglio compressione $V_{Rcd}$, calcolate come:
\begin{align}
    V_{Rd} &= \min \{V_{Rsd} ;  V_{Rcd} \} \label{eq:vrd}\\
    V_{Rsd} &= 0.9\,d\,\frac{A_{sw}}{s}\,f_{yd}\,(\cot\alpha + \cot\vartheta) \, \sin\alpha \\
    V_{Rcd} &= 0.9\,d\,b_w\,\alpha_c\,\nu\,f_{cd}\,\frac{\cot\alpha + \cot\vartheta}{1 + \cot^2\vartheta} 
\end{align}

Da queste equazioni le incognite libere sono l'inclinazione $\vartheta$ del puntone e il passo delle staffe $s$. 
Gli altri parametri sono tutti fissati da normativa o scelti preliminarmente.
Si sono scelte infatti staffe a due bracci di diametro $\varnothing_{st} = \SI{8}{\milli\metre}$.

Eguagliando la resistenza taglio compressione $V_{Rcd}$ al taglio sollecitante $V_{Ed}$ è possibile ricavare l'angolo $\vartheta$, da cui:
\begin{equation}
    \vartheta_\textup{design} = \max \left\{\frac{1}{2} \arcsin \frac{2\,V_{Ed}}{0.9\,d\,b_w\,\alpha_c\,\nu\,f_{cd}}  \; ; \; \SI{21.8}{\degree}\right\} \quad ,
\end{equation}
sostituendo tale valore dell'angolo nell'equazione del taglio trazione $V_{Rsd}$ e uguagliando anch' essa a $V_{Ed}$, si ottiene il passo massimo di progetto:
\begin{equation}
    s_\textup{design} = 0.9\,d\,A_{sw}\,f_{yd}\, \frac{\cot \vartheta_\textup{design}}{V_{Ed}} \quad .
\end{equation}

Tali valori di progetto del passo massimo sono riportati in tabella \ref{tab:V_passo}. 
Successivamente si è scelto il passo reale delle staffe con il quale si è proceduto poi alla verifica.
Si è cercato di avere le meno differenze di diverso passo possibile, in modo da non creare confusione in fase costruttiva dell'opera.


\begin{table}[htb]
    \centering
    \scriptsize
    \caption{ULS taglio design passi}
    \label{tab:V_passo}
    \begin{tabular}{
                        l
                        S[table-format=3.3]
                        S[table-format=2.1]
                        S[table-format=2.1]
                        S[table-format=3.1]
                }
    \toprule
    \multirow{2}{*}{Sez.} & {$V_{Ed}$} & {$\vartheta_\textup{design}$} & {$\vartheta_\textup{design, reale}$} & {$s_\textup{design}$} \\
    & {\si{[\kilo\newton]}} & {\si{[\degree]}} &  {\si{[\degree]}} &  {\si{[\milli\metre]}} \\

    \midrule
    A1 sx & 0.000 & 0.0 & 21.8 & $\infty$ \\
    A1 dx & 144.220 & 9.6 & 21.8 & 282.3 \\
    A2 sx & 236.080 & 16.2 & 21.8 & 172.5 \\
    A2 dx & 275.030 & 19.4 & 21.8 & 148.0 \\
    A3 sx & 275.730 & 19.4 & 21.8 & 147.7 \\
    A3 dx & 263.720 & 18.4 & 21.8 & 154.4 \\
    A4 sx & 254.150 & 17.6 & 21.8 & 160.2 \\
    A4 dx & 236.150 & 16.2 & 21.8 & 172.4 \\
    A5 sx & 255.130 & 17.7 & 21.8 & 159.6 \\
    A5 dx & 235.450 & 16.2 & 21.8 & 172.9 \\
    A6 sx & 227.870 & 15.6 & 21.8 & 178.7 \\
    A6 dx & 184.800 & 12.4 & 21.8 & 220.3 \\
    A7 sx & 149.290 & 9.9 & 21.8 & 272.7 \\
    A7 dx & 0.000 & 0.0 & 21.8 & $\infty$ \\
    \bottomrule
    \end{tabular}
    \end{table}
    

\paragraph{Verifica armature}
Fissato ora il passo reale $s$ delle staffe, eguagliando $V_{Rcd}$ con $V_{Rsd}$ è possibile calcolare il valore di theta in funzione di tutti gli altri parametri:
\begin{equation}
    \cot\vartheta_\textup{sezione} = \sqrt{-1 + \frac{\alpha_c \cdot b_w \cdot f_{cd} \cdot \nu \cdot s}{A_{sw}\cdot f_{yd} \cdot \sin\alpha}}
\end{equation}

In base al valore dell'angolo $\vartheta$ (o analogamente della sua cotangente) si attivano i tre diversi meccanismi resistenti e di conseguenza va scelto il valore di $\vartheta$ reale.
\begin{equation}
    \vartheta_\textup{reale} = 
\begin{cases}
    \vartheta_\textup{sezione} & \text{se} \quad \SI{21.8}{\degree}<\vartheta_\textup{sezione}<\SI{45}{\degree} \\
    \SI{21.8}{\degree} & \text{se} \quad \vartheta_\textup{sezione}\leq\SI{21.8}{\degree} \\
    \SI{45}{\degree} & \text{se} \quad \vartheta_\textup{sezione}\geq\SI{45}{\degree} \\
\end{cases}\quad,
\end{equation} 
dove nel primo caso i valori della resistenza a taglio lato calcestruzzo e lato acciaio saranno i medesimi, nel secondo la crisi avviene lato acciaio, mentre nell'ultimo lato calcestruzzo.

Con tali valori di $\vartheta_\textup{reale}$ è possibile infine verificare la sezione con le formule dei due meccanismi.

Nella tabella \ref{tab:V_with_armor} viene riportato tale procedimento per tutte le sezioni di verifica e la verifica risulta pertanto soddifìsfatta.
\begin{table}[htb]
    \centering
    \scriptsize
    \caption{ULS taglio con armatura}
    \label{tab:V_with_armor}
    \begin{tabular}{
        l
        c
        S[table-format=3.3]
        S[table-format=2.1]
        S[table-format=2.1]
        S[table-format=3.3]
        S[table-format=3.3]
        S[table-format=3.3]
        r
        S[table-format=2.1]}
    \toprule
    \multirow{2}{*}{Sez.} & \multirow{2}{*}{$A_{sw}$} & {$V_{Ed}$} & {$\vartheta_\textup{sezione}$} & {$\vartheta_\textup{reale}$} & {$V_{Rsd}$} & {$V_{Rcd}$} & {$V_{Rd}$} & \multicolumn{2}{l}{$V_{Ed}<V_{Rd}$} \\
    &&{\si{[\kilo\newton]}}&{\si{[\degree]}}&{\si{[\degree]}}&{\si{[\kilo\newton]}}&{\si{[\kilo\newton]}}&{\si{[\kilo\newton]}}&&{\si{[\percent]}}\\
    \midrule
    A1 sx & 2br.Ø8/220 & 0.000   & 16.9 & 21.8 & 185.076 & 303.346 & 185.076 & \checked & 0.0 \\
    A1 dx & 2br.Ø8/220 & 144.220 & 16.9 & 21.8 & 185.076 & 303.346 & 185.076 & \checked & 77.9 \\
    A2 sx & 2br.Ø8/140 & 236.080 & 21.3 & 21.8 & 290.834 & 303.346 & 290.834 & \checked & 81.2 \\
    A2 dx & 2br.Ø8/140 & 275.030 & 21.3 & 21.8 & 290.834 & 303.346 & 290.834 & \checked & 94.6 \\
    A3 sx & 2br.Ø8/140 & 275.730 & 21.3 & 21.8 & 290.834 & 303.346 & 290.834 & \checked & 94.8 \\
    A3 dx & 2br.Ø8/140 & 263.720 & 21.3 & 21.8 & 290.834 & 303.346 & 290.834 & \checked & 90.7 \\
    A4 sx & 2br.Ø8/140 & 254.150 & 21.3 & 21.8 & 290.834 & 303.346 & 290.834 & \checked & 87.4 \\
    A4 dx & 2br.Ø8/140 & 236.150 & 21.3 & 21.8 & 290.834 & 303.346 & 290.834 & \checked & 81.2 \\
    A5 sx & 2br.Ø8/140 & 255.130 & 21.3 & 21.8 & 290.834 & 303.346 & 290.834 & \checked & 87.7 \\
    A5 dx & 2br.Ø8/140 & 235.450 & 21.3 & 21.8 & 290.834 & 303.346 & 290.834 & \checked & 81.0 \\
    A6 sx & 2br.Ø8/160 & 227.870 & 19.9 & 21.8 & 254.480 & 303.346 & 254.480 & \checked & 89.5 \\
    A6 dx & 2br.Ø8/160 & 184.800 & 19.9 & 21.8 & 254.480 & 303.346 & 254.480 & \checked & 72.6 \\
    A7 sx & 2br.Ø8/220 & 149.290 & 16.9 & 21.8 & 185.076 & 303.346 & 185.076 & \checked & 80.7 \\
    A7 dx & 2br.Ø8/220 & 0.000   & 16.9 & 21.8 & 185.076 & 303.346 & 185.076 & \checked & 0.0 \\
    \bottomrule
    \end{tabular}
    \end{table}

\section{Zone di campata}
Nelle zone centrali delle campate si utilizza invece l'armatura minima precedentemente calcolata.
Avendo scelto staffe a due bracci Ø8, l'area vale $A_{st} = 2 \frac{\pi \varnothing^2}{4} = \SI{100.531}{\milli\metre\squared}$, corrispondente ad un numero di staffe:
\[
    n_{st, min} = \frac{A_{st, min}}{A_{st}} = \SI{4.476}{staffe\per\metre}    
\]
ed ad un passo di:
\[
    s_{st, max} = \frac{\SI{1000}{\milli\metre}}{n_{st,min}} = \SI{223.4}{\milli\metre} \implies \SI{220}{\milli\metre}  \quad .
\]

Con tale valore di 2br.Ø8/220 si ottiene il medesimo taglio resistente della sezione (ad esempio) A1 dx precedentemente calcolata in tabella \ref{tab:V_with_armor}, ovvero $V_{Rd} = \SI{185.076}{\kilo\newton}$.


\chapter{Trave: SLE}
\section{Verifica tensioni massime di esercizio}
In questo paragrafo si affronta la verifica delle tensioni massime di esercizio per le sezioni della trave già precedentemente definite agli SLU, pertanto le caratteristiche geometriche e di armature sono le medesime.

Le ipotesi alla base di tale verifica riguardano la capacità di resistenza del calcestruzzo: esso infatti non si considera resistente a trazione, ovvero è in stato 2.
Si ipotizza inoltre il comportamento lineare ed elastico di calcestruzzo e acciaio e con un coefficiente di omogenizzazione $n = 15$.

In riferimento alle \norma{NTC18} al capitolo \normaref{4.1.2.2.5} si deve avere che le tensioni di esercizio siano minori di quelle resistenti.
Si farà riferimento alla combinazione rara e a quella quasi permanente, con le tensioni resistenti riportate all'inizio di questo documento, al capitolo \ref{chap:materiali}.

Infine le formule che seguiranno faranno riferimento alla nomenclatura visibile in figura \ref{fig:sle_convenzione_noTraz}.

\begin{figure}[htb]
    \centering
    \includegraphics[height=0.25\textheight]{IMG/IPE_sle_cls_noTraz.pdf}
    \caption{Convenzione e nomenclatura utilizzata agli SLE, con ipotesi di CLS non resistente a trazione}
    \label{fig:sle_convenzione_noTraz}
  \end{figure}
  
Come primo passo occorre calcolare la posizione dell'asse neutro $nn$ con la coordinata $x$. 
Per farlo è possibile risolvere il sistema di equilibrio alla traslazione e momento, oppure, in alternativa, essendo le azioni assiali nulle, si ha che il momento statico rispetto l'asse neutro è nullo. 
Da quest'ultima equazione è perciò possibile determinare $x$.
\begin{align}
    0 &= S_{nn}  \\
    0 &= \frac{1}{2} b\,x^2 + n\,A_s^\prime\,(x - d^{\prime\prime}) - n\,A_S(d-x) \\
    &\quad\hookrightarrow x = \frac{n\,(A_s^\prime + A_s)}{b} \left( -1 + \sqrt{1 + \frac{2\,b\,(A_s^\prime \, d^{\prime\prime} + A_s\,d)}{n\,(A_s^\prime + A_s)^2}} \, \right)
\end{align}
Trovata la posizione dell'asse neutro la tensione sul calcestruzzo la si ricava grazie all'equazione di Bernouli-Navier:
\begin{equation}
    \sigma_c = \frac{M_{Ed}}{I_{nn}} \, x
\end{equation}
dove $I_{nn}$ è riferito alla sezione reagente, pertanto si considera solo la zona compressa di calcestruzzo in aggiunta alle armature opportunamente omogeneizzate:
\begin{equation}
    I_{nn} = \frac{b\, x^3}{12} + n\,A_s^\prime\,(x - d^{\prime\prime})^2 - n\,A_s(d-x)^2 \quad.
\end{equation}
Infine le tensioni sull'acciaio si calcolano grazie alla linearità del diagramma delle tensioni:
\begin{align}
    \sigma_s &= \frac{n\,\sigma_c}{x} \, (d-x) \\
    \sigma_s^\prime &= \frac{n\,\sigma_c}{x} \,(x - d^{\prime\prime})
\end{align}

Tale procedimento è stato applicato per ogni sezione e i risultati sono visibili nelle tabelle \ref{tab:SLE_verifica_CHAR} e \ref{tab:SLE_verifica_QP}, rispettivamente per la combinazione rara e per quella quasi permanente.
Come si può notare la verifica è soddisfatta in tutte le sezioni.
\begin{table}[H]
    \centering
    \scriptsize
    \caption{SLS CHAR con $\sigma_c = \SI{15.00}{\mega\pascal}$ e $\sigma_s = \sigma_s^\prime = \SI{360.00}{\mega\pascal}$}
    \label{tab:SLE_verifica_CHAR}
    \begin{tabular}{
        l
        c
        c
        S[table-format=3.3]
        S[table-format=3.2]
        S[table-format=2.2]
        S[table-format=3.2]
        S[table-format=3.2]
        c
        c
        c}
    \toprule
    \multirow{2}{*}{Sez.} & \multirow{2}{*}{$A_s$} & \multirow{2}{*}{$A_s^\prime$} & {$M_{Ed}$} 				             & {$x$}                 & {$\sigma_c$}          & {$\sigma_s$}          & {$\sigma_s^\prime$}   & \multicolumn{1}{c}{\multirow{2}{*}{$\sigma_c<\sigma_{cR}$}}& \multicolumn{1}{c}{\multirow{2}{*}{$\sigma_s<\sigma_{sR}$}}& \multicolumn{1}{c}{\multirow{2}{*}{$\sigma_s^\prime<\sigma_{sR}^\prime$}} \\


                        &                        &                               & {\si{[\kilo\newton\metre]}} 			& {\si{[\milli\metre]}} & {\si{[\mega\pascal]}} & {\si{[\mega\pascal]}} & {\si{[\mega\pascal]}} &                                                   &                                                     &  \\
    \midrule
    A1 inv & 2Ø18 & 2Ø18 & 0.000   & 116.55 & 0.00  & 0.00   & 0.00   & \checked & \checked & \checked \\
    C1     & 3Ø18 & 2Ø18 & 52.840  & 139.36 & 4.83  & 166.84 & 51.70  & \checked & \checked & \checked \\
    A2 inv & 6Ø18 & 2Ø18 & 127.100 & 185.66 & 9.34  & 207.09 & 109.96 & \checked & \checked & \checked \\
    C2     & 4Ø18 & 2Ø18 & 93.640  & 157.51 & 7.79  & 224.43 & 87.18  & \checked & \checked & \checked \\
    A3 inv & 6Ø18 & 2Ø18 & 136.290 & 185.66 & 10.02 & 222.06 & 117.91 & \checked & \checked & \checked \\
    C3     & 3Ø18 & 2Ø18 & 65.970  & 139.36 & 6.04  & 208.30 & 64.55  & \checked & \checked & \checked \\
    A4 inv & 6Ø18 & 2Ø18 & 123.070 & 185.66 & 9.05  & 200.52 & 106.47 & \checked & \checked & \checked \\
    C4     & 4Ø18 & 2Ø18 & 81.210  & 157.51 & 6.76  & 194.64 & 75.61  & \checked & \checked & \checked \\
    A5 inv & 6Ø18 & 2Ø18 & 166.360 & 185.66 & 12.23 & 271.06 & 143.93 & \checked & \checked & \checked \\
    C5     & 4Ø18 & 2Ø18 & 103.770 & 157.51 & 8.63  & 248.71 & 96.61  & \checked & \checked & \checked \\
    A6 inv & 6Ø18 & 2Ø18 & 135.560 & 185.66 & 9.97  & 220.88 & 117.28 & \checked & \checked & \checked \\
    C6     & 2Ø18 & 2Ø18 & 27.760  & 116.55 & 2.93  & 129.60 & 28.88  & \checked & \checked & \checked \\
    A7 inv & 3Ø18 & 2Ø18 & 53.400  & 139.36 & 4.89  & 168.61 & 52.25  & \checked & \checked & \checked \\
    \midrule
    A1     & 2Ø18 & 2Ø18 & 0.000   & 116.55 & 0.00  & 0.00   & 0.00   & \checked & \checked & \checked \\
    C1 inv & 2Ø18 & 3Ø18 & 0.000   & 111.06 & 0.00  & 0.00   & 0.00   & \checked & \checked & \checked \\
    A2     & 2Ø18 & 6Ø18 & 0.000   & 97.91  & 0.00  & 0.00   & 0.00   & \checked & \checked & \checked \\
    C2 inv & 2Ø18 & 4Ø18 & 0.000   & 106.17 & 0.00  & 0.00   & 0.00   & \checked & \checked & \checked \\
    A3     & 2Ø18 & 6Ø18 & 0.000   & 97.91  & 0.00  & 0.00   & 0.00   & \checked & \checked & \checked \\
    C3 inv & 2Ø18 & 3Ø18 & 0.000   & 111.06 & 0.00  & 0.00   & 0.00   & \checked & \checked & \checked \\
    A4     & 2Ø18 & 6Ø18 & 0.000   & 97.91  & 0.00  & 0.00   & 0.00   & \checked & \checked & \checked \\
    C4 inv & 2Ø18 & 4Ø18 & 0.000   & 106.17 & 0.00  & 0.00   & 0.00   & \checked & \checked & \checked \\
    A5     & 2Ø18 & 6Ø18 & 0.000   & 97.91  & 0.00  & 0.00   & 0.00   & \checked & \checked & \checked \\
    C5 inv & 2Ø18 & 4Ø18 & 0.000   & 106.17 & 0.00  & 0.00   & 0.00   & \checked & \checked & \checked \\
    A6     & 2Ø18 & 6Ø18 & 0.000   & 97.91  & 0.00  & 0.00   & 0.00   & \checked & \checked & \checked \\
    C6 inv & 2Ø18 & 2Ø18 & 0.000   & 116.55 & 0.00  & 0.00   & 0.00   & \checked & \checked & \checked \\
    A7     & 2Ø18 & 3Ø18 & 0.000   & 111.06 & 0.00  & 0.00   & 0.00   & \checked & \checked & \checked \\
    \bottomrule
    \end{tabular}
    \end{table}
%%%%%%%%%%%%%%%%%%%%%%%%%%%%%%%%%%%%%%%%%%%%%%%%%%%%%%%%%%%%%%%%%
\begin{table}[H]
    \centering
    \scriptsize
    \caption{SLS QP con $\sigma_c = \SI{11.25}{\mega\pascal}$ e $\sigma_s = \sigma_s^\prime = \SI{360.00}{\mega\pascal}$}
    \label{tab:SLE_verifica_QP}
    \begin{tabular}{
        l
        c
        c
        S[table-format=3.3]
        S[table-format=3.2]
        S[table-format=2.2]
        S[table-format=3.2]
        S[table-format=3.2]
        c
        c
        c}
    \toprule
    \multirow{2}{*}{Sez.} & \multirow{2}{*}{$A_s$} & \multirow{2}{*}{$A_s^\prime$} & {$M_{Ed}$}                      & {$x$}                 & {$\sigma_c$}          & {$\sigma_s$}          & {$\sigma_s^\prime$}   & \multicolumn{1}{c}{\multirow{2}{*}{$\sigma_c<\sigma_{cR}$}}& \multicolumn{1}{c}{\multirow{2}{*}{$\sigma_s<\sigma_{sR}$}}& \multicolumn{1}{c}{\multirow{2}{*}{$\sigma_s^\prime<\sigma_{sR}^\prime$}} \\


                        &                        &                               & {\si{[\kilo\newton\metre]}}      & {\si{[\milli\metre]}} & {\si{[\mega\pascal]}} & {\si{[\mega\pascal]}} & {\si{[\mega\pascal]}} &                                                   &                                                     &  \\
    \midrule
    A1 inv & 2Ø18 & 2Ø18 & 0.000   & 116.55 & 0.00 & 0.00   & 0.00   & \checked & \checked & \checked \\
    C1     & 3Ø18 & 2Ø18 & 32.310  & 139.36 & 2.96 & 102.02 & 31.61  & \checked & \checked & \checked \\
    A2 inv & 6Ø18 & 2Ø18 & 89.860  & 185.66 & 6.61 & 146.41 & 77.74  & \checked & \checked & \checked \\
    C2     & 4Ø18 & 2Ø18 & 62.110  & 157.51 & 5.17 & 148.86 & 57.83  & \checked & \checked & \checked \\
    A3 inv & 6Ø18 & 2Ø18 & 94.390  & 185.66 & 6.94 & 153.79 & 81.66  & \checked & \checked & \checked \\
    C3     & 3Ø18 & 2Ø18 & 36.930  & 139.36 & 3.38 & 116.60 & 36.13  & \checked & \checked & \checked \\
    A4 inv & 6Ø18 & 2Ø18 & 84.390  & 185.66 & 6.20 & 137.50 & 73.01  & \checked & \checked & \checked \\
    C4     & 4Ø18 & 2Ø18 & 52.860  & 157.51 & 4.40 & 126.69 & 49.21  & \checked & \checked & \checked \\
    A5 inv & 6Ø18 & 2Ø18 & 128.520 & 185.66 & 9.45 & 209.40 & 111.19 & \checked & \checked & \checked \\
    C5     & 4Ø18 & 2Ø18 & 79.040  & 157.51 & 6.58 & 189.43 & 73.59  & \checked & \checked & \checked \\
    A6 inv & 6Ø18 & 2Ø18 & 105.910 & 185.66 & 7.79 & 172.57 & 91.63  & \checked & \checked & \checked \\
    C6     & 2Ø18 & 2Ø18 & 19.710  & 116.55 & 2.08 & 92.02  & 20.51  & \checked & \checked & \checked \\
    A7 inv & 3Ø18 & 2Ø18 & 33.500  & 139.36 & 3.06 & 105.77 & 32.78  & \checked & \checked & \checked \\
    \midrule
    A1     & 2Ø18 & 2Ø18 & 0.000   & 116.55 & 0.00 & 0.00   & 0.00   & \checked & \checked & \checked \\
    C1 inv & 2Ø18 & 3Ø18 & 0.000   & 111.06 & 0.00 & 0.00   & 0.00   & \checked & \checked & \checked \\
    A2     & 2Ø18 & 6Ø18 & 0.000   & 97.91  & 0.00 & 0.00   & 0.00   & \checked & \checked & \checked \\
    C2 inv & 2Ø18 & 4Ø18 & 0.000   & 106.17 & 0.00 & 0.00   & 0.00   & \checked & \checked & \checked \\
    A3     & 2Ø18 & 6Ø18 & 0.000   & 97.91  & 0.00 & 0.00   & 0.00   & \checked & \checked & \checked \\
    C3 inv & 2Ø18 & 3Ø18 & 0.000   & 111.06 & 0.00 & 0.00   & 0.00   & \checked & \checked & \checked \\
    A4     & 2Ø18 & 6Ø18 & 0.000   & 97.91  & 0.00 & 0.00   & 0.00   & \checked & \checked & \checked \\
    C4 inv & 2Ø18 & 4Ø18 & 0.000   & 106.17 & 0.00 & 0.00   & 0.00   & \checked & \checked & \checked \\
    A5     & 2Ø18 & 6Ø18 & 0.000   & 97.91  & 0.00 & 0.00   & 0.00   & \checked & \checked & \checked \\
    C5 inv & 2Ø18 & 4Ø18 & 0.000   & 106.17 & 0.00 & 0.00   & 0.00   & \checked & \checked & \checked \\
    A6     & 2Ø18 & 6Ø18 & 0.000   & 97.91  & 0.00 & 0.00   & 0.00   & \checked & \checked & \checked \\
    C6 inv & 2Ø18 & 2Ø18 & 0.000   & 116.55 & 0.00 & 0.00   & 0.00   & \checked & \checked & \checked \\
    A7     & 2Ø18 & 3Ø18 & 0.000   & 111.06 & 0.00 & 0.00   & 0.00   & \checked & \checked & \checked \\
    \bottomrule
    \end{tabular}
\end{table}

\section{Verifiche di fessurazione}
In questo paragrafo vengono illustrate le verifiche di fessurazione per gli stati limite di esercizio in combinazione frequente e quasi caratteristica della sezione più sollecitata a momento positivo, ovvero la sezione di campata C5, come ben visibile nei diagrammi precedentemente riportati a pagine \pageref{fig:trave_SLS_FREQ_momento} e \pageref{fig:trave_SLS_QP_momento}.
Si riportano qui i valori delle sollecitazioni che verranno prese in esame:
\begin{equation}
    M_{Ed}^\textup{Freq.} = \SI{102.61}{\kilo\newton\metre} \quad
    M_{Ed}^\textup{QP} = \SI{79.04}{\kilo\newton\metre}
\end{equation}

\paragraph{Momento di  prima fessurazione}
Occorre come prima cosa controllare se la sezione si fessuri: in tal caso si procederà al calcolo della sua ampiezza.
Per farlo si calcola il momento di prima fessurazione 
\begin{equation}
    M_{cr} = \frac{ \sigma_{ct} }{ { h - x } } \, I_{nn} 
\end{equation}
dove 
\[
    \sigma_{ct} = \frac{ f_\textup{ctm} }{ 1.2 }  = \frac{ \SI{2.56}{\mega\pascal} }{ 1.2 } = \SI{2.14}{\mega\pascal}    
\]
mentre la distanza $x$ dall'asse neutro (calcolata ponendo il momento statico pari a zero come visto precedentemente) e il momento di inerzia $I_{nn}$ sono calcolati considerando anche il contributo a trazione del calcestruzzo. 
Si veda la nomenclatura adottata in figura \ref{fig:sle_convenzioneTraz}
\begin{figure}[htb]
    \centering
    \includegraphics[height=0.25\textheight]{IMG/IPE_sle_cls_Traz.pdf}
    \caption{Convenzione e nomenclatura utilizzata agli SLE, con ipotesi di CLS resistente a trazione}
    \label{fig:sle_convenzioneTraz}
  \end{figure}

Essi valgono:
\[
    \begin{split}
    x 
    &= \frac{ n \cdot \left( A_{s} \cdot d + A^{\prime}_{s} \cdot d^{\prime\prime} \right) + b \cdot \frac{ \left( h \right) ^{ 2 } }{ 2 } }{ n \cdot \left( A_{s} + A^{\prime}_{s} \right) + b \cdot h } \\
    &= \frac{ 15 \cdot \left( \SI{1017.85}{\milli\metre\squared} \cdot \SI{460.00}{\milli\metre} + \SI{508.92}{\milli\metre\squared} \cdot \SI{40.00}{\milli\metre} \right) + \SI{300.00}{\milli\metre} \cdot \frac{ \left( \SI{500.00}{\milli\metre} \right) ^{ 2 } }{ 2 } }{ 15 \cdot \left( \SI{1017.85}{\milli\metre\squared} + \SI{508.92}{\milli\metre\squared} \right) + \SI{300.00}{\milli\metre} \cdot \SI{500.00}{\milli\metre} } \\
    &= \SI{259.27}{\milli\metre}  \\
    %
    I_{nn} 
    &= b \cdot \frac{ \left( x \right) ^{ 3 } }{ 3 } + n \cdot A^{\prime}_{s} \cdot \left( x - d^{\prime\prime} \right) ^{ 2 } + n^{\prime} \cdot b \cdot \frac{ \left( h - x \right) ^{ 3 } }{ 3 } + n \cdot A_{s} \cdot \left( d - x \right) ^{ 2 } \\
    &= \SI{300.00}{\milli\metre} \cdot \frac{ \left( \SI{259.27}{\milli\metre} \right) ^{ 3 } }{ 3 } + 15 \cdot \SI{508.92}{\milli\metre\squared} \cdot \left( \SI{259.27}{\milli\metre} - \SI{40.00}{\milli\metre} \right) ^{ 2 } \\
    & \quad\quad + 1 \cdot \SI{300.00}{\milli\metre} \cdot \frac{ \left( \SI{500.00}{\milli\metre} - \SI{259.27}{\milli\metre} \right) ^{ 3 } }{ 3 } + 15 \cdot \SI{1017.85}{\milli\metre\squared} \cdot \left( \SI{460.00}{\milli\metre} - \SI{259.27}{\milli\metre} \right) ^{ 2 } \\
    &= \SI{4120094011.14}{\milli\metre\tothe{4}}  \\
    %
    M_{cr} 
    &= \frac{ \sigma_{ct} }{ n^{\prime} } \cdot \frac{ I_{nn} }{ h - x }  
    = \frac{ \SI{2.14}{\mega\pascal} }{ 1 } \cdot \frac{ \SI{4120094011.14}{\milli\metre\tothe{4}} }{ \SI{500.00}{\milli\metre} - \SI{259.27}{\milli\metre} } 
    = \SI{36.58}{\kilo\newton\metre}  < 
        \begin{bmatrix}
            M_{Ed}^\textup{Freq.} \\
            M_{Ed}^\textup{QP}
        \end{bmatrix} \quad \notchecked
    \end{split}
\]
Si può notare quindi come $M_{cr}$ sia minore di $M_{Ed}$ in entrambe le combinazioni, perciò la trave è fessurata. 
\subsection{Metodo esplicito}
Le \norma{NTC18} ai capitoli \normaref{4.1.2.2.4} e \normaref{C4.1.2.2.4.5} regolano il procedimento per il calcolo dell'ampiezza di fessurazione con un metodo esplicito, in cui si calcola il vero valore dell'ampiezza di fessura $w_k$ e lo si confronta con dei limiti di riferimento, e un metodo semplificato in cui si effettua il calcolo di $w_k$ attraverso una tabella.

I valori limite si basano sulle condizioni ambientali e di sensibilità alla corrosione dell'armatura.
Assumendo di avere delle condizioni ambientali ordinarie e delle armature poco sensibili alla corrosione si hanno i seguenti limiti di ampiezza:
\begin{equation}
    \begin{split}
        w_2 &= \SI{0.3}{\milli\metre} \quad \text{Comb. Frequente} \\ 
        w_3 &= \SI{0.4}{\milli\metre} \quad \text{Comb. Quasi permanente} \\
    \end{split}
\end{equation}

Il calcolo dell'ampiezza vera di fessurazione è invece dato da:
\begin{equation}
    w_k = 1.7\, \varepsilon_\textup{sm} \cdot \Delta_{sm}
\end{equation}
dove $\varepsilon_\textup{sm}$ rappresenta la deformazione media delle barre d'armatura, mentre $\Delta_{sm}$ rappresenta la distanza media tra le fessure.
\begin{equation}
    \varepsilon_\textup{sm} = \frac{ \sigma_{s} - k_{t} \cdot \frac{ f_\textup{ctm} }{ \rho_\textup{eff} } \cdot \left( 1 + \alpha_{e} \cdot \rho_\textup{eff} \right) }{ E_{s} } > 0.6 \cdot \frac{ \sigma_{s} }{ E_{s} } \quad ,
\end{equation}
in cui tutti i vari termini sono definiti in normativa nel seguente modo:
\[
\begin{split}
    \alpha_{e} &= \frac{ E_{s} }{ E_{cm} }  = \frac{ \SI{210000.00}{\mega\pascal} }{ \SI{31475.81}{\mega\pascal} } = 6.67  
    \\
    h_\textup{c,eff} 
    &= \min { \left( 2.5 \cdot \left( h - d \right) ,\  \frac{ h - x }{ 3 } ,\  h \cdot \frac{1} { 2 } \right) }  \\ 
    & = \min { \left( 2.5 \cdot \left( \SI{500.00}{\milli\metre} - \SI{460.00}{\milli\metre} \right) ,\  \frac{ \SI{500.00}{\milli\metre} - \SI{157.51}{\milli\metre} }{ 3 } ,\  \SI{500.00}{\milli\metre} \cdot \frac{1} { 2 } \right) } \\
    &= \SI{100.00}{\milli\metre}  \\
    A_\textup{c,eff} &= b \cdot h_\textup{c,eff}  = \SI{300.00}{\milli\metre} \cdot \SI{100.00}{\milli\metre} = \SI{30000.00}{\milli\metre\squared}  \\
    \rho_\textup{eff} &= \frac{ A_{s} }{ A_\textup{c,eff} }  = \frac{ \SI{1017.85}{\milli\metre\squared} }{ \SI{30000.00}{\milli\metre\squared} } = \num{3.39e-2 }  \\
\end{split}
\]
\begin{align*}
k_{1} &= 0.8 \; \text{barre ad aderenza migliorata}  & k_{2} &= 0.5 \; \text{nel caso della flessione} \\
k_{3} &= 3.4 \; \text{fisso} & k_{4} &= 0.425 \;\text{fisso} \\
c &= \SI{20.000}{\milli\metre} \; \text{ricoprimento barre} & k_t &= 
    \begin{cases}
        0.6 \quad \text{Frequente} \\
        0.4 \quad \text{Quasi permanente}
    \end{cases}
\end{align*}
\[
    \Delta_{sm} = \frac{ k_{3} \cdot c + k_{1} \cdot k_{2} \cdot k_{4} \cdot \frac{ \Phi }{ \rho_\textup{eff} } }{ 1.7 }  = \frac{ 3.4 \cdot \SI{20.000}{\milli\metre} + 0.8 \cdot 0.5 \cdot 0.425 \cdot \frac{ \SI{18.000}{\milli\metre} }{\num{3.393e-2 } } }{ 1.7 } = \SI{93.053}{\milli\metre}  
\]
\paragraph{Combinazione quasi permanente}
Dalla precedente tabella \ref{tab:SLE_verifica_QP} si ha che l'asse neutro e la tensione dell'acciaio valgono:
\[
    x = \SI{157.51}{\milli\metre} \quad \sigma_{s} = \SI{189.43}{\mega\pascal}
\]
È perciò ora possibile applicare le formule della normativa, ottenendo:
\[
\begin{split}
    \varepsilon_\textup{sm} 
    &= \frac{ \sigma_{s} - k_{t} \cdot \frac{ f_\textup{ctm} }{ \rho_\textup{eff} } \cdot \left( 1 + \alpha_{e} \cdot \rho_\textup{eff} \right) }{ E_{s} }  \\
    &= \frac{ \SI{189.43}{\mega\pascal} - 0.4 \cdot \frac{ \SI{2.56}{\mega\pascal} }{\num{3.39e-2 } } \cdot \left( 1 + 6.67 \cdot\num{3.39e-2 } \right) }{ \SI{210000.00}{\mega\pascal} } \\
    &=\num{7.25e-4 }  \\
    %
    \varepsilon_\textup{sm,lim} 
    &= 0.6 \cdot \frac{ \sigma_{s} }{ E_{s} }  
    = 0.6 \cdot \frac{ \SI{189.43}{\mega\pascal} }{ \SI{210000.00}{\mega\pascal} } =\num{5.41e-4 }  \\
    %
    \varepsilon_\textup{sm,vero} 
    &= \max { \left( \varepsilon_\textup{sm} ,\  \varepsilon_\textup{sm,lim} \right) }  
    = \max { \left(\num{7.25e-4 } ,\ \num{5.41e-4 } \right) } 
    =\num{7.25e-4 }  \\
    %
    w_{k} &=  1.7 \cdot \varepsilon_\textup{sm,vero} \cdot \Delta_{sm}  
    =  1.7 \cdot\num{7.255e-4 } \cdot \SI{93.053}{\milli\metre} 
    = \SI{0.115}{\milli\metre} < w_3 = \SI{0.4}{\milli\metre} \quad \checked\\
    \end{split}
    \]
    
 \paragraph{Combinazione frequente}
 Analogamente a quanto fatto sopra e guardando la tabella \ref{tab:SLE_verifica_CHAR}, si ha:
\[
    x = \SI{157.51}{\milli\metre} \quad \sigma_{s} = \SI{245.925}{\mega\pascal}
\] 
da cui:
\[
\begin{split}
    \varepsilon_\textup{sm} 
    &= \frac{ \sigma_{s} - k_{t} \cdot \frac{ f_\textup{ctm} }{ \rho_\textup{eff} } \cdot \left( 1 + \alpha_{e} \cdot \rho_\textup{eff} \right) }{ E_{s} }  \\
    &= \frac{ \SI{245.925}{\mega\pascal} - 0.6 \cdot \frac{ \SI{2.565}{\mega\pascal} }{\num{3.393e-2 } } \cdot \left( 1 + 6.672 \cdot\num{3.393e-2 } \right) }{ \SI{210000.000}{\mega\pascal} } \\
    &=\num{9.062e-4 }  \\
    %
    \varepsilon_\textup{sm,lim} 
    &= 0.6 \cdot \frac{ \sigma_{s} }{ E_{s} }  
    = 0.6 \cdot \frac{ \SI{245.925}{\mega\pascal} }{ \SI{210000.000}{\mega\pascal} } 
    =\num{7.026e-4 }  \\
    %
    \varepsilon_\textup{sm,vero} 
    &= \max { \left( \varepsilon_\textup{sm} ,\  \varepsilon_\textup{sm,lim} \right) }  
    = \max { \left(\num{9.062e-4 } ,\ \num{7.026e-4 } \right) } 
    =\num{9.062e-4 }  \\
    %
    w_{k} 
    &= 1.7 \cdot \varepsilon_\textup{sm,vero} \cdot \Delta_{sm}  
    =  1.7 \cdot\num{9.062e-4 } \cdot \SI{93.053}{\milli\metre} \ 
    = \SI{0.143}{\milli\metre}  < w_2 = \SI{0.3}{\milli\metre} \quad \checked\\
\end{split}
\]
\subsection{Metodo semplificato}
Il metodo semplificato consiste nel controllare che le armature adottate e la spaziature tra esse siano inferiori ai limiti imposti dalle \normaref{Tab. C4.1.II-III}.

Considerando la condizione più sfavorevole con $\sigma_{s} = \SI{245.925}{\mega\pascal} \simeq \SI{240}{\mega\pascal}$ il diametro massimo da adottare è un Ø20, mentre la spaziatura massima è \SI{250}{\milli\metre}. 
Avendo adottato 4Ø18, il diametro è inferiore a quello massimo e la spaziatura massima è sicuramente maggiore della spaziatura reale avendo una base di \SI{300}{\milli\metre}. 
Pertanto anche con il metodo semplificato la verifica a fessurazione è soddisfatta \checked.

\section{Verifiche di deformabilità}
\chapter{Pilastro}

In questo capitolo si progetterà e verificherà il pilastro P27 secondo l'azione assiale precedentemente riportata in figura \ref{fig:AndamentoP27} di pagina \pageref{fig:AndamentoP27}, di cui si riporta qui sotto l'azione massima alla base del pilastro:
\begin{equation}
    N_{Ed} = \SI{1416.00}{\kilo\newton} \quad N_{EQP} = \SI{829.35}{\kilo\newton}
\end{equation}
\section{Progetto pilastro alla base e verifica compresisone}
Per determinare l'area minima di calcestruzzo e di acciaio le \norma{NTC18} propongono le seguenti formule al \normaref{4.1.6.1.2}:
\begin{align}
    A_{s,min} &= 0.10\,\frac{N_{Ed}}{f_{yd}}  > 0.003 \, A_c \\
    A_c &= 0.90\,\frac{N_{Ed}}{f_{cd}}
\end{align}
i quali diventano rispettivamente $\SI{361.87}{\milli\metre\squared}$ e $\SI{299.93}{\milli\metre}$.

Scegliendo 2+2Ø12 = $\SI{452.16}{\milli\metre\squared}$ e  $\SI{300}{\milli\metre}$ come base del pilastro, la verifica risulta:
\begin{equation}
    N_{Rd} = 0.8 A_c f_{fcd} + A_s f_{yd} = \SI{1197}{\kilo\newton} < N_{Ed} = \SI{1416.00}{\kilo\newton} \quad \notchecked
\end{equation}

Viene pertanto aumentata la base a $\SI{350}{\milli\metre}$, che però ciò implica il dover disporre di una barra aggiuntiva, in quanto si avrebbe la distanza tra le barre di spigolo maggiore di $\SI{300}{\milli\metre}$.
Viene inoltre aumentato il diametro a 14, perciò si avranno 3+3Ø14 = $\SI{923}{\milli\metre\squared}$.
Occorre inoltre dover aumantare l'azione solecitante, perché calcolata con un predisionamento della precedente base. Vengono aggiunti quindi \SI{9.75}{\kilo\newton} al $G_1$, corrispondenti al peso proprio aggiuntivo, per l'intera altezza dell'edificio.
L'azione SLU diviene quindi $N_{Ed} = \SI{1425.75}{\kilo\newton}$. 
La verifica a compresisone risulta soddisfatta:

\begin{equation}
    N_{Rd} = 0.8 A_c f_{fcd} + A_s f_{yd} = \SI{1749.77}{\kilo\newton} > N_{Ed} = \SI{1425.75}{\kilo\newton} \quad \checked \quad .
\end{equation}
Il quantitavo massimo di armatura è inoltre soddisfatto:
\begin{equation}
    A_s =
      \SI{923}{\milli\metre\squared} 
    \begin{cases}
      < 0.04 \, A_c = 0.04 \, b \, h = \SI{9000}{\milli\metre\squared} \\
      > 0.003 \, A_c = 0.003 \, b \, h = \SI{367.50}{\milli\metre\squared}
    \end{cases}
    \quad \checked
  \end{equation}

Per quanto riguarda l'armatura trasversale a taglio si deve avere:
\begin{align}
    s_{st,min} &= \min\left( 12\varnothing_{long.} \, ; \, \SI{250}{\milli\metre} \right)= \SI{168}{\milli\metre}\\
    \varnothing_{st, min} &= \max\left( 0.25\varnothing_{long.} \, ; \, \SI{6}{\milli\metre} \right)= \SI{6}{\milli\metre}
\end{align}
Da cui si scelgono staffe a due bracci Ø8/\SI{150}{\milli\metre}

\section{Verifica presso-flessione}
Sebbene il pilastro sia sottoposto a puro sforzo assiale, la normativa richiede di considerare un minimo quantito di eccentricità calcolabile come
\[
    e_0 = \max \left(\frac{l_0}{200};  \SI{20}{\milli\metre} \right)  = \max \left(\frac{3500}{200};  20 \right)\si{\milli\metre} = \SI{20}{\milli\metre}
\]
la quale genera un momento flettente pari a 
\begin{equation}
    M_{Ed} = N_{Ed}\cdot e_0  = \SI{1425.75}{\kilo\newton} \cdot (\SI{20}{\milli\metre} = \SI{28.52}{\kilo\newton\metre}
\end{equation}

Per verificarne la resistenza a presso-flessione si fa uso del diagramma M-N che permette di rappresentare la resistenza della sezione  per tutti i valori possibili di azione assiale.
Tale diagramma è riportato in figura \ref{fig:MN_diagram_primo_ordine}, dimostrando che la verifica è soddifìsfatta.
\begin{figure}[H]
    \centering
    \includegraphics[width=\textwidth]{IMG/M_N_diagram_pilastro_primo_ordine.pdf}
    \caption{Diagramma MN presso-flesisone}
    \label{fig:MN_diagram_primo_ordine}
\end{figure}


\section{Verifica di stabilità}
Occorre innanzitutto controllare se la snellezza reale del pilastro superi o meno quella limite fissata da normativa. 
Nel caso venga superata, occorre effettuare una più accurata analisi degli effetti del secondo ordine.

La snellezza per il pilastro più alto, vale:
\begin{equation}
    \lambda = \frac{l_0}{i} = \frac{\SI{3.5}{\metre}}{\frac{\SI{0.350}{\metre}}{\sqrt{12}}} = \num{34.641} \, ,
\end{equation}
dove $l_0$ è la lunghezza libera d'inflessione, ed è pari all' altezza del pilastro, avendo considerato uno schema statico di appoggio-appoggio, 
e $i$ è il raggio giratore d'inerzia calcolato come $\sqrt{\frac{J}{A}}$ della sezione non fessurata.

La snellezza limite, secondo le \norma{NTC18} alla \normaref{[4.1.41]}, nel caso di pilastri singoli, può essere calcolata in via semplificata come:
\begin{equation}
    \lambda_\textup{limite} = \frac{25}{n} =  \frac{25}{\dfrac{N_{Ed}}{A_c f_{cd}}} = \num{27.582}
\end{equation}
in cui $n$ rappresenta l'azione assiale adimensionalizzata.

Essendo $\lambda > \lambda_\textup{limite}$ occore prendere in considerazione gli effetti del secondo ordine.
Per farlo verrà fatto uso del metodo basato sulla curvatura nominale come riportato al \normaref{\S 5.8.8} dell' \norma{EC2}.
Tale metodo prevede di calcolare il contributo del momento flettente di secondo ordine, da sommare al momento di primo ordine (che tiene conto delle imperfezioni), e di confrontarlo poi con il momento resistente della sezione, fissato lo sforzo assiale sollecitante. 
Ovvero, in formule:
\begin{equation}
    M_{Ed} = M_{0,Ed} + M_2 < M_{Rd}(N_{Ed})
\end{equation}

\paragraph{Momento del primo ordine}
L'eccentricità $e_0$ che tiene conto del non allineamento dei carichi è calcolabile come già visto in precedenza
\[
    e_0 = \max \left(\frac{l_0}{200};  \SI{20}{\milli\metre} \right)  = \max \left(\frac{3500}{200};  20 \right)\si{\milli\metre} = \SI{20}{\milli\metre}
\]
mentre l'eccentricità $e_a$ che tiene conto delle imperfezioni geometriche, dal \normaref{\S 5.2(7)}, è calcolabile come
\[
    e_a = \vartheta_i \cdot \frac{l_0}{2} = \SI{8.75}{\milli\metre}
\]
in cui si è adottato $\vartheta_i$ pari a $1/200$.

Pertanto il momento di primo ordine diviene
\begin{equation}
    \begin{split}
        M_{0,Ed} &= N_{Ed}\cdot(e_0 + e_a)  = \SI{1425.75}{\kilo\newton} \cdot (\SI{20}{\milli\metre} + \SI{8.75}{\milli\metre} ) = \SI{40.99}{\kilo\newton\metre} \\
        M_{0,EQP} &= N_{EQP}\cdot(e_0 + e_a)  = \SI{839.1}{\kilo\newton} \cdot (\SI{20}{\milli\metre} + \SI{8.75}{\milli\metre} ) = \SI{20.12}{\kilo\newton\metre}
    \end{split}
\end{equation}


\paragraph{Momento del secondo ordine}
L'eccentricità $e_2$ che tiene conto del contributo degli effetti del secondo ordine è calcolabile, secondo il \normaref{\S 5.8.8.2(3)} dell' \norma{EC2} come 
\begin{equation}
    e_2 = \frac{1}{r} \frac{l_0 ^ 2}{c}
\end{equation}
dove $\frac{1}{r}$ è la curvatura ed è calcolabile come mostrato sotto, mentre $c$ è un fattore che dipende dalla distribuzione della curvatura e  per sezioni trasversali costanti, può essere preso pari a $10$.

La curvatura $\frac{1}{r}$ è calcolabile tramite dei coefficienti correttivi a partire da quella nominale nel seguente modo:
\begin{equation}
    \frac{1}{r} = K_r \cdot K_\varphi \cdot \frac{1}{r_0}
\end{equation}
in cui $K_r$ dipende dal carico assiale, mentre $K_\varphi$ tiene conto della viscosità:
\[
\begin{split}
    K_r &= \frac{n_\textup{u} - n}{n_\textup{u} - n_\textup{bal}} = \frac{1.20825 - 0.821541}{1.20825 - 0.4} = 0.478452\\
    n_\textup{u} &= 1 + \omega = 1.20825 \\
    \omega &= \frac{A_\textup{s, tot} f_{yd}}{A_c f_{cd}} = 0.20825 \\
    K_\varphi &= \max \left(1 + \beta \cdot \varphi_\textup{ef}; 1 \right) = 1.267596	 \\
    \beta &= 0.35 + \frac{f_{ck}}{200} - \frac{\lambda}{150} = 0.24406 \\
    \varphi_\textup{ef} &= \varphi_{\infty, t_0, h_0} \frac{M_{0,EQP}}{M_{0,Ed}}= 1.863 \frac{24.12}{40.99} = 1.096 \\
        &\quad \text{nel quale $\varphi_{\infty, t_0, h_0}$ è stato preso dalla tabella \normaref{11.2.VI} delle \norma{NTC18}} \\
    \frac{1}{r_0} &= \frac{\varepsilon_{yd}}{0.45 d} = \frac{\frac{f_{yd}}{E_s}}{0.45 d} = \SI{1.336e-5}{\milli\metre\tothe{-1}}
\end{split}
\quad .
\] 
Perciò infine la curvatura vale 
\[
    \frac{1}{r} = K_r \cdot K_\varphi \cdot \frac{1}{r_0} = 0.153183 \cdot 1.267596	 \cdot \num{1.336e-5} = \SI{8.101e-6}{\milli\metre\tothe{-1}}
\]
e l'eccentricità del secondo ordine e il momento divengono
\begin{equation}
    \begin{split}
        e_2 &=\SI{8.101e-6}{\milli\metre\tothe{-1}} \cdot \frac{(\SI{3500}{\milli\metre})^2}{10} = \SI{9.924}{\milli\metre} \\
        M_2 &=  N_{Ed}\cdot e_2  = \SI{1425.75}{\kilo\newton} \cdot \SI{9.924}{\milli\metre} = \SI{14.15}{\kilo\newton\metre} 
    \end{split}
\end{equation}

\paragraph{Momento totale}
Sommando  i due contributi $M_{0,Ed}$ e $M_2$ si ottiene il momento solleccitante totale  $M_{Ed} = \SI{55.14}{\kilo\newton\metre}$.
È quindi possibile verificare tramite il diagramma M-N riportato in figura \ref{fig:MN_diagram_secondo_ordine} che è all'interno del dominio resistente $M_{Rd}$, e pertanto la verifica di stabilità risulta soddisfatta.

\begin{figure}[H]
    \centering
    \includegraphics[width=\textwidth]{IMG/M_N_diagram_pilastro_secondo_ordine.pdf}
    \caption{Diagramma MN con sollecitazione flettente tenente conto degli effetti del secondo ordine tramite il metodo della curvatura nominale}
    \label{fig:MN_diagram_secondo_ordine}
\end{figure}
\end{document} 